%% sections/11_conclusion.tex

\section{결론}
\label{sec:conclusion}

본 논문은 LLM 투기적 서빙에서 동시에 발생하는 두 가지 비효율을 분석하고,
이를 통합적으로 해결하는 알고리즘 \algname 을 제안하였다.

첫 번째 비효율은 \emph{워크로드 다양성}이다.
투기적 디코딩의 효과는 워크로드 유형, 모델 크기, 모델 계열에 따라
naive n-gram의 붕괴($-90\%$)에서 suffix의 강한 가속(mix $+232\%$)까지 극도로 다르며,
최적 방법이 (모델 $\times$ 조건)마다 달라 단일 설정으로 최적화가 원리적으로 불가능하다.
특히 code agent나 추론 특화 모델과 같은 최신 응용에서 이 문제는 더욱 심각하다.

두 번째 비효율은 \emph{CPU 유휴}이다.
첫 번째 문제를 해결하기 위해 이중 백엔드를 채택하면
CPU 이용률이 $\le5.6\%$에 머무는 구조적 슬랙이 필연적으로 발생한다.

\algname 은 이 두 문제를 인과적으로 연결된 하나의 구조로 해결한다.
워크로드$\times$모델 크기$\times$벤더 3축 게이팅이 첫 번째 문제를 해결하고,
그 구조가 낳은 CPU 슬랙을 NUMA 바인딩과 스텝 경계 정렬 branchy 보조 작업으로 수확한다.

DGX B200$\times$8 에서 8개 모델(Qwen$\cdot$Llama$\cdot$DeepSeek 계열)을
실제 서빙 trace(ShareGPT, WildChat, LMSYS-Chat, SWE-bench, MBPP, HumanEval)로
측정한 결과, spec-decode는 이질 트래픽에서 8개 모델 전체에 대해 순 양성
($+83\%$--$+232\%$)을 유지하면서도 최적 방법은 (모델, 조건)마다 달랐다.
동일 환경에서 실측한 클러스터 라우터 llm-d 가 최선 단일 방법을 능가한 경우는
18\%(10/56)에 그쳐, 경량 워크로드 게이팅이 대부분의 가치를 포착함을 확인하였다.
CPU 슬랙 하베스팅은 개발 플랫폼에서 NUMA 바인딩 $+3.13\%$(warm-only)의
robust 이득과 branchy 보조 작업 최대 $+5.28\%$의 예비 이득을 회수한다.

\algname 이 제시하는 핵심 원칙은 단순하다:
\emph{게이팅이 슬랙을 낳고, 슬랙을 수확하라.}
서버 내 유휴 자원은 고정 비용이 아니라 수확 가능한 기회이다.

\section*{감사의 글}

본 논문의 작성 과정에서 Anthropic 사의 Claude
(모델 ID: \texttt{claude-sonnet-4-6})를
측정 자료 분석, 초고 작성 보조, \LaTeX\ 편집에 활용하였다.
알고리즘 설계의 모든 결정, 실험 계획, 측정 결과 해석,
학술적 주장의 정확성에 대한 책임은 저자에게 있다.
